\title{Parameter-Efficient Orthogonal Finetuning via Factorization}

\begin{abstract}
As large foundation models gain prominence, training them anew from scratch becomes exceedingly resource-intensive. Consequently, finding effective strategies for adapting these powerful models to specific downstream applications is of increasing importance. In this work, we focus on Orthogonal Finetuning (OFT)—a structured approach for adapting foundation models to downstream tasks. While OFT achieves strong generalization performance, its reliance on orthogonal matrices entails a substantial number of trainable parameters, attributed to their high dimensionality. To mitigate this limitation, we analyze OFT through the lens of information transmission and highlight several critical criteria for improving parameter efficiency. Drawing inspiration from the Cooley-Tukey fast Fourier transform algorithm—which excels in efficient data transformation—we propose a compact orthogonal parameterization based on butterfly matrix structures. Integrating this parameterization into OFT yields a new, parameter-efficient finetuning technique: Orthogonal Butterfly (BOFT). BOFT generalizes OFT, treating the latter as a specific instance within its broader framework. We validate BOFT with extensive empirical studies, demonstrating effective adaptation of large vision transformers, large language models, and text-to-image diffusion models to a range of tasks spanning both vision and language domains.
\end{abstract}

\section{Introduction}

Breakthrough models like ChatGPT~\cite{brown2020language,chen2021evaluating} and Stable Diffusion~\cite{rombach2022high} have underscored the extraordinary generalization capacities of large foundation models. However, the swift escalation in their performance has coincided with explosive growth in parameter counts 

(\emph{e.g.}, GPT-3 contains approximately 175 billion parameters), presenting steep challenges for practitioners aiming to train such models independently. As a result, advancing the field depends crucially on methods for leveraging these pre-trained models without requiring full retraining. Thus, efficient techniques to adapt existing foundation models to specific tasks have become central. Three main strategies are prevalent for this purpose: \emph{model finetuning}~\cite{hulora2022,zaken2022bitfit,guo2020parameter,raffel2020exploring,qiu2023controlling,zhang2022adaptive,chavan2023one}, which leaves the overall architecture intact but selectively updates parts of the model parameters; \emph{adapter tuning}~\cite{rebuffi2017learning,houlsby2019parameter,pfeiffer2020adapterfusion,lin2020exploring,he2021towards}, wherein new trainable modules are appended to the original network; and \emph{prompt tuning}~\cite{li2021prefix,lester2021power}, which involves introducing trainable prefix tokens into the input sequence. Of these, model finetuning stands out due to its simplicity, strong empirical performance, and critically, the absence of added inference-time overhead.

The core idea underlying model finetuning is to maintain a certain level of similarity between the pretrained model and its finetuned counterpart according to specific metrics, thus retaining the knowledge acquired during pretraining. Typically, existing finetuning strategies employ a reduced learning rate—a heuristic that helps minimize the divergence between the pretrained and finetuned models. Considering the inherent structure of weight matrices, an alternative, more theoretically grounded approach seeks to maintain the relational characteristics of these matrices, particularly the pairwise angles among neurons. This consideration gives rise to the Orthogonal Finetuning~(OFT)~\cite{qiu2023controlling} framework, wherein all neurons within a given layer are subjected to transformation by a shared orthogonal matrix, thereby guaranteeing the conservation of inter-neuronal angles throughout the finetuning procedure. Although OFT has achieved notable improvements in generalizability and convergence, especially when applied to text-to-image diffusion models~\cite{qiu2023controlling}, it is burdened by the substantial number of parameters stemming from the high-dimensional nature of orthogonal matrices. To alleviate this issue, OFT employs a block-diagonal orthogonal matrix, thereby reducing parameter count. Yet, this gain in parameter efficiency introduces its own set of challenges: the imposed fixed sparsity pattern and the fact that the orthogonal transformations are executed independently across different blocks. While the block-diagonal structure does limit parameter usage, it also risks imposing problematic inductive biases (for instance, diminishing representational flexibility or precluding the approximation of classical linear mappings), and, crucially, there is still no established method for determining an optimal sparsity pattern.

A central challenge in this context is to produce a dense orthogonal matrix that remains parameter-efficient. At first, this task may appear unattainable, as a fully dense orthogonal matrix of dimension $d$ typically involves $\mathcal{O}(d^2)$ parameters. However, our approach diverges from conventional methods by constructing a dense orthogonal matrix through the multiplication of several sparse, factorized matrices. This method stems from the observation that parameter count can be minimized by accepting increased computational demands in exchange for reduced space complexity. Because the dense orthogonal matrix is encoded as a product of sparse matrices, multiple sequential multiplications must occur during every training iteration. To conceptualize matrix factorization, we frame the creation of a dense orthogonal matrix as analogous to transmitting information through a network. Within this analogy, forming a dense orthogonal matrix by multiplying sparse matrices can be likened to information flow across a grid-shaped graph. This perspective provides a systematic way to devise various strategies for sparse matrix factorization, allowing us to cap the number of trainable parameters while maintaining enough expressiveness to approximate dense matrices.

In order to further enhance parameter efficiency within the information transmission paradigm, we turn to butterfly networks as utilized in the Cooley-Tukey fast Fourier transform algorithm~\cite{cooley1965algorithm}, which are noted for their ability to exchange information rapidly among different nodes~\cite{klasing1994broadcasting}. The butterfly architecture underpinning the Cooley-Tukey algorithm introduces a natural blueprint for a sparse factorization method, known as butterfly factorization. Through this method, a $d$-dimensional dense matrix can be composed as the product of $\mathcal{O}(\log d)$ sparse matrices, each possessing only $\mathcal{O}(d)$ nonzero elements. By imposing orthogonality constraints on each constituent sparse matrix, we derive a resource-efficient orthogonal parameterization utilizing just $\mathcal{O}(d\log d)$ parameters—a substantial saving over standard parameter counts. Exploiting this compressed orthogonal parameterization, we introduce a new finetuning technique, termed Orthogonal Butterfly~(BOFT). BOFT generalizes and contains OFT as a particular scenario, creating a broad orthogonal finetuning framework. Both the block-diagonal configuration and butterfly structure share the key attribute of sparsity, integrating data sparsity into orthogonal matrices to curtail the trainable parameter count. Notably, it is informative to juxtapose our methodology with the low-rank design utilized in LoRA~\cite{hulora2022}; both low-rank and sparse matrices represent prominent categories within structured matrices~\cite{candes2011robust}, each attaining parameter efficiency through their respective structures.

In contrast to the block-diagonal configuration adopted in OFT to balance expressiveness and regularization, BOFT leverages the butterfly structure to create a more continuous spectrum between matrices from the complete orthogonal group (ensuring expressivity) and the identity matrix (ensuring regularity). This design choice allows the discovery of a more constrained hypothesis space within the orthogonal group for specific downstream applications. Given the prevalent role of the butterfly structure in accelerating various linear transformations—such as discrete Fourier and discrete cosine transforms—we posit that employing this structured parameterization imparts a beneficial inductive bias, which not only enhances generalization but also guards against overfitting. Our rationale is that the butterfly structure can readily reconstruct a wide array of classical linear operators, an ability absent from the block-diagonal framework utilized in OFT. The principal contributions of this work are summarized as follows:

\begin{itemize}[leftmargin=*,nosep]
    \item We investigate parameter efficiency in the context of orthogonal finetuning, introducing a novel information transmission perspective. Within this perspective, constructing a dense, parameter-efficient orthogonal matrix is reformulated as an information flow problem on a grid-like graphical model.
    \item Drawing inspiration from butterfly arrangements exemplified in the Cooley-Tukey algorithm, we introduce Orthogonal Butterfly, a new parameter-efficient strategy for orthogonal finetuning situated within our information transmission framework.
    \item We furnish several theoretical arguments explaining how BOFT is able to greatly reduce the number of learnable parameters while maintaining ample expressive capacity and stable training dynamics. The structure enables matrix factorization, which additionally supports an appealing form of weight interpolation (illustrated in Figure~\ref{fig:boft_interpolation}).
    \item For the first time, we extend orthogonal finetuning to domains beyond controllable text-to-image synthesis~\cite{qiu2023controlling}, highlighting its promise as a broadly applicable finetuning mechanism. We apply BOFT across a variety of adaptation challenges spanning computer vision and natural language processing. Empirical results indicate that BOFT consistently surpasses current state-of-the-art alternatives, underscoring both its high parameter efficiency and robust generalization.
\end{itemize}

\textbf{Parameter-efficient finetuning (PEFT)}. With the continual growth in size and capability of foundation models (\emph{e.g.}, \cite{brown2020language,radford2021learning,rombach2022high,kirillov2023segment}), there is heightened attention to adapting these models to specific downstream tasks in a way that minimizes the number of parameters updated~\cite{treviso2023efficient,ding2023parameter,lialin2023scaling}. A variety of PEFT techniques have been proposed~\cite{houlsby2019parameter,aghajanyan2020intrinsic,hulora2022,edalati2022krona,wang2022adamix,gheini2021cross,zaken2022bitfit,guo2020parameter,sung2021training,ansell2022composable,lester2021power,li2021prefix,vu2022spot,he2021towards,mao2021unipelt,karimi2021compacter,liu2022few,sung2022lst,chen2022parameter,jia2022visual,chen2022adaptformer,zhang2022neural,jie2023fact,lian2022scaling,luo2023towards}, among which reparameterization-driven approaches~\cite{aghajanyan2020intrinsic,hulora2022,edalati2022krona} are most aligned with our study. LoRA~\cite{hulora2022} modifies pretrained weight matrices through the addition of products of two low-rank matrices, demonstrating notable success on various natural language processing benchmarks. Whereas LoRA fixes the rank of these matrices uniformly, alternative strategies \cite{zhang2022adaptive,valipour2022dylora,zhang2023increlora} have been developed to allow the rank to vary across layers, thus ensuring a balanced distribution of parameter resources. The straightforwardness of employing a low-rank reparameterization for weights has led to its widespread use~\cite{dettmers2023qlora,zi2023delta,chavan2023one}. Drawing on insights from hyperspherical energy’s role in understanding generalization~\cite{liu2018learning,liu2021orthogonal}, orthogonal finetuning (OFT)~\cite{qiu2023controlling} offers a contrasting and effective paradigm, especially for tuning text-to-image diffusion models. Here, OFT introduces a learnable orthogonal matrix to transform neurons within each layer, yielding improved generalization and more robust, stable learning when compared to LoRA. Nevertheless, OFT typically incurs a higher count of trainable parameters than LoRA. Thus, increasing the parameter-efficiency of OFT emerges as a promising direction. Additionally, whether OFT extends readily to adaptation scenarios beyond its original application to text-to-image diffusion models~\cite{qiu2023controlling} remains uncertain. BOFT addresses this efficiency limitation by incorporating butterfly factorization, thereby enabling broader demonstrations of the effectiveness of orthogonal finetuning in a variety of adaptation tasks.

\textbf{Butterfly structures}. The radix-2 Cooley-Tukey algorithm transforms the computation of the $N$-point discrete Fourier transform into smaller subproblems of size $\frac{N}{2}$, recursively constructing a butterfly structure that decomposes into a product of several sparse matrices (with this product referred to as a butterfly matrix). Such butterfly matrices have proven valuable for orthogonal matrix parameterizations—circumventing the need for pivoting during Gaussian elimination and facilitating computational efficiency~\cite{parker1995random}—as well as stabilizing recurrent neural network training~\cite{jing2017tunable} and enabling efficient kernel approximations~\cite{munkhoeva2018quadrature}. The use of butterfly-based parameterizations for learning efficient linear transforms has been explored by \cite{dao2019learning,dao2019kaleidoscope}, and their utility in efficient neural network training has been further studied by \cite{chen2022pixelated,dao2022monarch}. Beyond these, butterfly structures contribute to accelerating matrix-vector products~\cite{michielssen1996multilevel,o2010algorithm}, providing compact matrix approximations~\cite{li2015butterfly}, and optimizing network data transmission~\cite{klasing1994broadcasting,li2003linear,rajasingh2016transmission}. In contrast to prior literature, the present work centers on leveraging butterfly structures specifically to improve the parameter efficiency of OFT.

\section{An Information Transmission View on Orthogonal Finetuning}

We begin by introducing the foundational concepts of OFT. In order to finetune a pretrained weight matrix, OFT reformulates the updated weight matrix as the product of a trainable orthogonal matrix and the fixed, pretrained weights. Unlike LoRA, which adds a low-rank matrix to adjust the weights, OFT modifies them through multiplication by an orthogonal matrix. LoRA attains parameter-efficiency via low-rank approximation, whereas the original OFT achieves efficiency by employing a block-diagonal orthogonal matrix structure~\cite{qiu2023controlling}; reducing the block size enhances parameter efficiency. A visual comparison is provided in Figure~\ref{fig:comp_lora}. The rationale behind orthogonal transformation in finetuning is to maintain neuron pairwise angles~\cite{liu2018learning,liu2021orthogonal,qiu2023controlling}, thereby preserving much of the semantic information learned during pretraining. Specifically, for a pretrained linear layer $\bm{W}^0\in\mathbb{R}^{d\times n}$, OFT learns an orthogonal matrix $\bm{R}\in\mathbb{R}^{d\times d}$, adjusting the forward computation from $\bm{z}=(\bm{W}^0)^\top\bm{x}$ to $\bm{z}=(\bm{R}\bm{W}^0)^\top\bm{x}$, where $\bm{x}\in\mathbb{R}^{d}$ and $\bm{z}\in\mathbb{R}^{n}$ denote the input and output vectors, respectively. To guarantee zero initialization, the orthogonal matrix $\bm{R}$ is initially set to the identity. Preserving orthogonality of $\bm{R}$ during finetuning leverages the Cayley parameterization, following \cite{qiu2023controlling,liu2021orthogonal}, with $\bm{R}=(\bm{I}+\bm{Q})(\bm{I}-\bm{Q})^{-1}$, where $\bm{Q}$ is a skew-symmetric matrix satisfying $\bm{Q}=-\bm{Q}^\top$. For enhanced parameter efficiency, the orthogonal matrix $\bm{R}$ is structured as block-diagonal, that is, $\text{diag}(\bm{R}_1,\bm{R}_2,\cdots,\bm{R}_r)$, where each block $\bm{R}_i\in\mathbb{R}^{b\times b}$ is itself an orthogonal matrix and $br = d$. However, this approach assumes that each neuron’s dimensions (that is, the entries of a column of $\bm{W}^0$) are partitioned into $r$ groups, with each group subjected to a distinct orthogonal transformation, which may not align with the underlying neuron structure. While this block-diagonal design demonstrates strong empirical results, there is little rationale for segmenting neuron dimensions based merely on index, suggesting a preference for dense orthogonal matrices. This leads to a key inquiry: \emph{Is it possible to devise a dense orthogonal matrix that retains parameter efficiency?}

To tackle this issue, we introduce an approach that represents a dense orthogonal matrix as a product of several sparse orthogonal matrices. To systematically examine the sparsity configurations in orthogonal matrix decompositions, we conceptualize the construction of a dense orthogonal matrix within OFT as analogous to an information transmission scenario. In particular, forming a dense matrix $\bm{R}\in\mathbb{R}^{d\times d}$ as a product of $m$ square matrices, namely $\thickmuskip=2mu \medmuskip=2mu \bm{R}=\bm{B}_m\bm{B}_{m-1}\cdots\bm{B}_1$, is equivalent to modeling the flow of information across a $d\times (m+1)$ grid of nodes, as depicted in Figure~\ref{fig:info_trans}. This perspective draws on the intuition that a dense $d$-dimensional square matrix encapsulates complete connectivity between $d$ source and $d$ target nodes. For the matrix $\bm{R}$, a zero in the $R_{ij}$ entry signifies that information cannot be conveyed from node $j$ to node $i$, whereas a nonzero value means transmission is possible. Hence, expressing the dense matrix $\bm{R}$ as a sequence of factors $\bm{B}_m\bm{B}_{m-1}\cdots\bm{B}_1$ can be interpreted as progressive information propagation guided by the structure of each factor’s corresponding graph. The transmission sequence starts with $\bm{B}_1$ and culminates at $\bm{B}_n$. As illustrated in Figure~\ref{fig:info_trans}, we use the example of the decomposition $\bm{R}=\bm{B}_5\bm{B}_4\bm{B}_3\bm{B}_2\bm{B}_1$, with its associated sparsity patterns and graph representation shown. The illustrated graph represents an unrolled view of the matrix product process, where $\bm{B}_i$ defines the connectivity from the nodes at level $i$ to those at level $(i+1)$. Explicitly, the $(j_1, j_2)$ entry in $\bm{B}_i$ corresponds to the existence of a directed link from the $j_2$-th node on level $i$ to the $j_1$-th node on level $i+1$, with zero indicating the absence of such an edge. To ensure that the product $\bm{B}_5\bm{B}_4\bm{B}_3\bm{B}_2\bm{B}_1$ results in a dense matrix, it is necessary that all nodes from the first level are capable of reaching any node on the sixth level through information transmission. Examining, for instance, $\bm{R}=\bm{B}_4\bm{B}_3\bm{B}_2\bm{B}_1$, which connects sources at the first level to recipients at the fifth, we observe that node 1 fails to transmit information to node 3, and so the sparsity pattern of the product includes a zero at position $(3,1)$. A similar relationship between the induced graphs and sparsity patterns can be observed for the products $\bm{R}=\bm{B}_3\bm{B}_2\bm{B}_1$ and $\bm{R}=\bm{B}_2\bm{B}_1$. In a more general setting, achieving a dense matrix $\bm{R}\in\mathbb{R}^{d\times d}$ requires the factor matrices $\bm{B}_m,\cdots,\bm{B}_1$ to together generate a network of directed edges over a $d\times (m+1)$ grid such that each edge connects nodes only between adjacent columns (levels), ensuring that information from every node in the initial layer can reach all nodes at the final, $(m+1)$-th, layer.

Adopting the information transmission perspective offers valuable insights into the sparsity characteristics of factorization matrices in OFT. As depicted in Figure~\ref{fig:blk_diag}, the original OFT presents a block-diagonal structure in $\bm{R}$. Although this configuration decreases the quantity of trainable parameters, it fails to yield a fully dense matrix $\bm{R}$. The objective is to synthesize a dense orthogonal matrix from $m$ sparse orthogonal matrices while utilizing the fewest possible effective trainable parameters. Within the information transmission paradigm, two primary criteria emerge: \emph{(i)} \textbf{dense connectivity}, where each node in the initial level possesses at least one pathway to every node in the terminal level; and \emph{(ii)} \textbf{minimum free edges}, ensuring that the aggregate number of edges is minimized subject to orthogonality constraints. Orthogonality introduces nuanced restrictions on edge configuration between consecutive levels. Specifically, for any matrix $\bm{B}_i$ to attain full rank (which is essential for orthogonality), there must be $d$ edges to establish a bijective correspondence between all nodes at the $i$-th level and those at the $(i=1)$-th level, requiring at least $d$ inter-level edges (\emph{e.g.}, in Figure~\ref{fig:info_trans}, this count is 4). These $d$ edges, being indispensable for orthogonality, should be excluded from the total edge count, as they are fixed and not trainable (\emph{e.g.}, in a $d\times d$ orthogonal matrix with $d$ non-zero components, their values must be $\pm 1$). The orthogonality condition, which reduces the number of required parameters in comparison to a dense matrix, imposes further dependencies among the edges. For instance, in a $2\times 2$ orthogonal matrix, a zero at position $(1,1)$ enforces a zero at position $(2,2)$ (\emph{i.e.}, omitting one edge necessitates the omission of another). Consequently, depending on the chosen pattern of edge connections, orthogonality may lead to additional edge insertions or removals. When analyzing the non-zero structure of a composed orthogonal matrix, the information transmission approach proves particularly apt for OFT, since the primary concern lies with the locations of the non-zero, trainable entries in $\bm{R}$, while the actual numerical values are inconsequential.

A straightforward dense linkage between two hierarchical levels comprises $\mathcal{O}(d^2)$ connections (that is, represented by a single dense orthogonal matrix), resulting in $d^2-d$ total edges (for $d=4$, this equates to 12 edges). Figure~\ref{fig:info_trans} presents an example of a viable matrix decomposition that utilizes just $10$ edges overall—fewer than those in a complete dense orthogonal matrix. This approach allows us to analyze the parameter-efficiency of OFT graphically, and makes the derivation of valid factorizations straightforward within this scheme. Motivated by a distinctive connectivity structure from the Cooley-Tukey algorithm, known as butterfly graphs~\cite{cooley1965algorithm}, we exploit their ability to connect $d$ input nodes to $d$ output nodes using only $\mathcal{O}(d\log d)$ edges, achieving high connectivity with fewer parameters. For instance, whereas the topology depicted in Figure~\ref{fig:info_trans} achieves dense connections using $10$ edges, the butterfly network architecture accomplishes this with just $8$ edges. In the following, we describe how incorporating the butterfly topology can further improve parameter efficiency.

\section{Orthogonal Parameterization by Butterfly Factorization}

Initially developed as part of the Cooley-Tukey algorithm for the fast Fourier transform, the butterfly structure enables rapid computation by promoting information spread—analogous to our own scenario where each initial node can relay information to every terminal node. A change in the frequency domain under the Fourier transform induces widespread alterations in the spatial domain, mirroring the complete information transmission objective here. Furthermore, butterfly networks are widely adopted in computer network design~\cite{solihin2015fundamentals,li2011linear} to facilitate effective data dissemination. Letting $k\geq 2$ be a power of two, we define the \emph{butterfly factor} $\bm{B}^F(k)$ as follows:

\begin{equation}\label{eq:but_factor}
\begin{aligned}
    \bm{B}^F(k)=\begin{bmatrix}
    \text{diag}(\bm{d}_1) & \text{diag}(\bm{d}_2) \\
    \text{diag}(\bm{d}_3) & \text{diag}(\bm{d}_4)
    \end{bmatrix}\in\mathbb{R}^{k\times k},
\end{aligned}
\end{equation}

where each $\bm{d}_i\in\mathbb{R}^{\frac{k}{2}}$ for $i=1,2,3,4$ are specified vectors. Building on this, for $d=2^N$, we recursively construct the $d$-dimensional \emph{butterfly matrix} $\bm{B}(d)\in\mathbb{R}^{d\times d}$ as

\begin{equation}
\begin{aligned}
    \!\!\bm{B}(d)=\tilde{\bm{B}}(d,d)\!\cdot\!\begin{bmatrix}
    \bm{B}_1(\frac{d}{2})\!\!\!\!\! & \bm{0} \\
    \bm{0}\!\!\!\!\! &\bm{B}_2(\frac{d}{2})
    \end{bmatrix}= \tilde{\bm{B}}(d,d)\tilde{\bm{B}}(d,\frac{d}{2})\cdots\tilde{\bm{B}}(d,2),
\end{aligned}
\end{equation}

where $\bm{B}_1(\frac{d}{2})$ and $\bm{B}_2(\frac{d}{2})$ denote two butterfly matrices, each of dimension $\frac{d}{2}$. The \emph{butterfly component} is formulated as $\thickmuskip=2mu \medmuskip=2mu \tilde{\bm{B}}(d,k)=\text{diag}(\bm{B}^F_1(k),\cdots,\bm{B}^F_{d/k}(k))$, representing a $d \times d$ block-diagonal matrix whose block size is $k$. Here, each $\bm{B}^F_i(k)$ (for every $i$) refers to a butterfly factor introduced in Equation~\ref{eq:but_factor}. With this foundation, the butterfly matrix is employed for parameterizing an orthogonal matrix. To do this, it suffices that every block-diagonal factor $\tilde{\bm{B}}(d,k)$, for all $k$, in the decomposition of $\bm{B}(d)$, is itself orthogonal. To illustrate, we examine $\tilde{\bm{B}}(d,2)$, the block-diagonal case with blocks of size $2$. Ensuring orthogonality of each $2 \times 2$ block is straightforward—by leveraging the Cayley transform or $2$-dimensional rotations for parameterization—as also adopted by \cite{liu2021orthogonal,qiu2023controlling}. The sparsity pattern of any butterfly component is simply a permutation of the nonzero structure in $\tilde{\bm{B}}(d,2)$, allowing all butterfly components to be efficiently realized as orthogonal matrices. This provides an efficient scheme for orthogonal matrix parameterization rooted in compositions of $2 \times 2$ orthogonal blocks. Extending the framework following \cite{chen2022pixelated}, we introduce the block butterfly component $\tilde{\bm{B}}^b(d,k)$, where every entry in $\bm{d}_i$ for any $i$ is promoted to a $b\times b$ matrix. In order to maintain orthogonality for $\tilde{\bm{B}}^b(d,2)$, each $2b \times 2b$ block matrix is accordingly parameterized to be orthogonal. For larger blocks ($k>2$), the butterfly components $\tilde{\bm{B}}^b(d,k)$ are structured by permuting the blockwise nonzero entries of $\tilde{\bm{B}}^b(d,2)$, which likewise enables them to be parameterized as orthogonal matrices. By composing these components, the forward operation in BOFT is expressed as

\begin{equation}
    \begin{aligned}\nonumber
    \bm{z}=\big(\bm{R}(m,b)\cdot\bm{W}^0\big)^\top\bm{x},~~~~\text{s.t.}~\left\{\bm{R}(m,b)=\prod_{i=1}^m \tilde{\bm{B}}^b_{(d,i)}~~\&~~\underbrace{\big(\tilde{\bm{B}}^b_{(d,j)}\big)^\top\tilde{\bm{B}}^b_{(d,j)}=\tilde{\bm{B}}^b_{(d,j)}\big(\tilde{\bm{B}}^b_{(d,j)}\big)^\top\!=\bm{I}_d}_{\forall j\in [1, m]}\right\},
    \end{aligned}
\end{equation}

Here, we abbreviate $\tilde{\bm{B}}^b(d,2^{m-i+1})$ as $\tilde{\bm{B}}^b_{(d,i)}$ for ease of notation, with $\bm{I}_d$ representing the identity matrix in $\mathbb{R}^{d\times d}$. The orthogonal matrix $\bm{R}(m,b)\in\mathbb{R}^{d\times d}$ is built from a sequence of orthogonal butterfly matrices. For compactness, BOFT parameterized as $\bm{R}(m, \frac{b}{2})$ is denoted as BOFT($m$, $b$) where $b\geq 2$. Notably, when $\thickmuskip=2mu \medmuskip=2mu m=1$, BOFT($1$, $b$) is equivalent to the block-diagonal OFT~\cite{qiu2023controlling} having block size $b$. Likewise, choosing BOFT($1$, $d$) returns the standard OFT~\cite{qiu2023controlling} that utilizes a full, unconstrained orthogonal matrix. Using BOFT($\log \frac{2d}{b}$, $b$), the block butterfly structure $\bm{B}^{\frac{b}{2}}(d)$ serves as $\bm{R}$, resulting in a dense orthogonal matrix. Generally, BOFT($m$, $b$) employs a total of $\frac{1}{2}(b-1)dm$ trainable parameters for finetuning a linear transformation of dimensions $d\times n$. When employing a butterfly configuration, i.e., $m=\log d, b=2$, the total number of parameters becomes $\mathcal{O}(d\log d)$. In contrast, the full dense orthogonal OFT requires $\mathcal{O}(d^2)$ parameters, and the block-diagonal OFT with $r$ blocks uses $\mathcal{O}(bd)$. Consequently, achieving a dense orthogonal matrix in original OFT necessitates setting $b=d$ for the block size, whereas BOFT offers the flexibility to select any $b$ to the same end.

\textbf{BOFT Identity Initialization}. Finetuning strategies typically initialize from the weights of the original pretrained network to avoid significant divergence from the base model. For instance, LoRA initializes its low-rank matrices to zero. In BOFT, we adopt an initialization where all butterfly components start as identity matrices, which amounts to using zero initialization for the associated skew-symmetric matrices in the Cayley transform.

\textbf{BOFT Multiplicative Dropout}. LoRA~\cite{hulora2022} employs a Dropout mechanism for its low-rank weights to mitigate overfitting. While standard Dropout~\cite{srivastava2014dropout} fits naturally in LoRA, this is not the case for BOFT, as its weight update is multiplicative. To resolve this issue, we introduce a multiplicative Dropout specifically for BOFT. Since $\bm{R}(m,b)$ consists of $m$ orthogonal butterfly matrices, which themselves can be reorganized into $2b \times 2b$ block-diagonal orthogonal matrices, the multiplicative Dropout approach entails randomly selecting a fraction $p_1$ of the butterfly matrices and within each, a proportion $p_2$ of their diagonal blocks, subsequently resetting these chosen blocks to the identity.

\section{Intriguing Insights and Discussions}

\textbf{Expressivity of BOFT}. The butterfly structure combined with permutations is capable of exactly representing various well-known fast linear transforms~\cite{dao2019learning,dao2019kaleidoscope} (\emph{e.g.}, fast Fourier transform, Hadamard transform). However, the capacity of our orthogonal butterfly matrix to closely approximate arbitrary orthogonal matrices has not been thoroughly explored. To investigate this, we conducted an empirical study in which a random, dense orthogonal matrix~\cite{anderson1987generation} of size $512\times 512$ was approximated. Figure~\ref{fig:para_eff} presents results averaged over 10 different initializations. Here, the y-axis represents the approximation error while the x-axis shows the number of effective trainable parameters. Each curve of a fixed color corresponds to BOFT with a specific block size; the point furthest to the left marks the performance of BOFT($1$,$b$) (i.e., the original block-diagonal OFT of block size $b$). In general, BOFT achieves greater parameter efficiency compared to OFT. For instance, BOFT($9$,$2$) demonstrates higher expressivity than BOFT($1$,$16$), yet uses substantially fewer parameters. Choosing a smaller $b$ in combination with a larger $m$ tends to yield more parameter-efficient configurations for BOFT. As an example, BOFT($6$,$4$) utilizes significantly fewer parameters but attains similar approximation errors as BOFT($2$,$16$). Fundamentally, the butterfly matrix forms a more constrained and structured subset within the orthogonal group (relative to block-diagonal constructions), thereby making BOFT provably more expressive than OFT when the block size is held constant.

\begin{theorem}[Expressivity of BOFT]\label{thm:exp_boft}
    BOFT possesses greater expressive power than OFT under an identical block size. To fully approximate the set of orthogonal matrices of size $d$, one can employ products of butterfly matrices as in $\bm{B}_{d-1,1}(d)\bm{B}^\top_{d-1,2}(d)\cdots\bm{B}_{1,1}(d)\bm{B}^\top_{1,2}(d)$, where each $\bm{B}_{i,j}(d)$ denotes a butterfly matrix.
\end{theorem}

Based on Theorem~\ref{thm:exp_boft}, a natural extension for BOFT is suggested—namely, the final orthogonal transformation can be generalized to the form $\thickmuskip=2mu \medmuskip=2mu \bm{R}^G(m_1,b_1,m_2,b_2,l)=\bm{R}_{l,1}(m_1,b_1)\bm{R}_{l,2}^T(m_2,b_2)\cdots\bm{R}_{1,1}(m_1,b_1)\bm{R}_{1,2}^T(m_2,b_2)$, with $\bm{R}_{i,j}^T(m,b)$ indicating the specific orthogonal matrices used within BOFT. When setting $m_1=m_2=\log d$, $b_1=b_2=2$, and $l=d-1$, $\bm{R}^G(m_1,b_1,m_2,b_2,l)$ has the capacity to express every member of the orthogonal group. This form of matrix product is also referred to as the kaleidoscope hierarchy~\cite{dao2019kaleidoscope}. Nevertheless, it is important to recognize that higher expressivity does not unconditionally result in superior fine-tuning outcomes; indeed, complete fine-tuning, despite its comprehensive expressiveness, frequently does not achieve satisfactory results. Therefore, balancing expressivity and regularization emerges as a central issue for the robustness of fine-tuned models. BOFT enhances the trainable parameter space via the imposition of structural priors, enabling a more optimal balance between expressivity and regularity.

\textbf{Spectral properties}. Orthogonal finetuning typically results in more favorable spectral characteristics than LoRA, as it fully maintains the spectral norm of the original pretrained weight matrix $\bm{W}^0$. This preservation can be demonstrated using the singular value decomposition: $\bm{W}^0 = \bm{U}\bm{\Sigma}\bm{V}^\top$, where $\bm{U}$ and $\bm{V}$ are orthogonal matrices and $\bm{\Sigma}$ is diagonal containing the singular values. In both OFT and BOFT, an orthogonal matrix $\bm{R}$ is applied to the left, resulting in the updated weight matrix $\bm{R}\bm{U}\bm{\Sigma}\bm{V}^\top$. This operation leaves the maximum singular value (i.e., the spectral norm of $\bm{W}^0$) unchanged. Maintaining the spectral norm in this way has been shown to enhance both stability during optimization and generalization ability~\cite{miyato2018spectral,yoshida2017spectral}. Additional mathematical insights into this phenomenon can be found in Appendix~\ref{app:sn_property}.

\textbf{Orthogonal finetuning as learning bilinear similarity}. The BOFT approach can also be interpreted in terms of learning a bilinear similarity measure, namely $\bm{w}^0_i\bm{R}\bm{x}$, where $\bm{w}^0_i$ denotes the $i$-th column (neuron) of $\bm{W}^0$. Here, BOFT can be framed as optimizing the bilinear similarity matrix $\bm{R}$ under an orthogonality constraint, establishing an intrinsic link with the frameworks of distance metric learning~\cite{xing2002distance} and bilinear forms~\cite{roman2005advanced}.

\textbf{Inductive bias and generalization in BOFT}. In BOFT, $\bm{R}(m,b)$ typically characterizes a structured segment of the orthogonal group, thus restricting the hypothesis space. This structural constraint acts as an inductive bias, and we propose that the particular inductive bias introduced by butterfly factorization is advantageous for generalization. This is supported by its patterned resemblance to widely used linear transforms~\cite{dao2019kaleidoscope}, including the discrete Fourier, sine/cosine, and Hadamard transforms. Furthermore, the sparsity present in the matrix factorization within BOFT could offer further implicit inductive biases~\cite{gunasekar2017implicit,li2018algorithmic,liu2019neural}.

\textbf{Comparison to butterfly-based sparse training}. Several studies~\cite{dao2019learning,dao2019kaleidoscope,chen2022pixelated,dao2022monarch} have explored sparse training utilizing the butterfly parameterization. These works generally concentrate on applying the butterfly structure directly to reparameterize weight matrices, training neural networks anew. In particular, \cite{dao2022monarch} investigates fine-tuning by initially projecting pretrained weights onto a modified butterfly matrix, subsequently optimizing the projected elements for specific downstream applications. In contrast, BOFT introduces a substantially different fine-tuning paradigm, employing transformations of the weights through layer-shared matrices.

\input{rebuttal-results}

\section{Concluding Remarks and Limitations}

In this work, we introduce Orthogonal Butterfly, a broadly applicable parameter-efficient fine-tuning technique for foundation models grounded in the butterfly architecture. Our primary innovation in enhancing parameter efficiency is to construct a dense orthogonal matrix as a product of several sparse orthogonal matrices. To facilitate the discovery of viable matrix decompositions, we develop a graphical framework based on information transmission. Through this lens, we identify that the butterfly structure can efficiently accomplish our objective of decomposing sparse orthogonal matrices. We present experimental results demonstrating BOFT’s effectiveness in fine-tuning models across domains, including large language models, substantial vision architectures, and text-to-image generative models. The empirical evidence substantiates BOFT’s capacity as a versatile fine-tuning approach.

Nevertheless, BOFT does come with certain drawbacks. Since the final orthogonal transformation within BOFT is obtained by multiplying multiple orthogonal matrices, the associated training time overhead slightly exceeds that of OFT. Addressing this training runtime inefficiency remains a subject for future exploration. Notably, once fine-tuning is complete, the orthogonal matrices derived via BOFT can be directly fused into the pretrained network, ensuring there is no added latency during inference. Furthermore, it remains uncertain whether the butterfly network constitutes the optimal design for information transmission. Our information transmission framework also opens avenues for cross-disciplinary innovation, drawing from the field of computer networking, which rigorously investigates network topology efficiency. This suggests the potential to discover even more effective network structures for constructing dense orthogonal matrices.

\end{document}